\documentclass[11pt,letterpaper]{article}
\usepackage[top=2.0cm, bottom=3cm, left=2.0cm, right=2.0cm]{geometry}
\usepackage[utf8]{inputenc}
\usepackage[T1]{fontenc}
\usepackage[spanish]{varioref}
\usepackage[activeacute, spanish, es-tabla]{babel}
\usepackage{fancyhdr}
\usepackage{multicol}
\usepackage{float}
\usepackage{textcomp}
\usepackage{ae,aecompl}
\usepackage{amssymb,amsmath}
\usepackage[pdftex]{graphicx}
\pagestyle{fancy} 
\pagenumbering{gobble}
\renewcommand{\headrulewidth}{0pt} 
\setlength{\headsep}{20pt} 
\setlength{\headheight}{65pt} 
\setlength{\textheight}{600pt} 
\setlength{\columnsep}{15pt} 
\newcommand{\universidad}{\small{Universidad Técnica Federico Santa María}}
\newcommand{\campus}{\small{Jorge Contreras 201573547-6}}
\newcommand{\semestre}{\small{Juan Pablo Jorquera  201573533-6}}

% Definiciones de Título e Integrantes de Experiencia
\newcommand{\titulo}{Informe Caso - Tarea 2 BD}
\newcommand{\integrantes}{\begin{tabular}{c}

\end{tabular}}

\renewcommand{\maketitle}
{
\thispagestyle{fancy}
\begin{center}
\begin{Large}
\textbf{\titulo}\\
\end{Large}
\end{center}
\vspace{0.3cm}
}


%ENCABEZADO

\fancyhead[R]{\begin{minipage}[b]{0.405\textwidth}
\begin{center}
\universidad \\ 
\campus \\ 
\lab \\ 
\semestre
\end{center}
\end{minipage}}
\fancyhead[L]{\vspace{15pt}\includegraphics[height=1.6cm]{Escudo.png}}
%%%%%%%%%%%%%%%%%%%%%%%%%%%%%%%%%%%%%%%%%%%%%%%%
%                                              %
% AQUI TERMINAN LAS DEFINICIONES DE ENCABEZADO %
% Y EMPIEZA EL CUERPO DEL DOCUMENTO            %
%                                              %
%%%%%%%%%%%%%%%%%%%%%%%%%%%%%%%%%%%%%%%%%%%%%%%%

\begin{document}
\maketitle
\section{Modelo Conceptual}
\vspace{0.6cm}
\begin{center}
	\includegraphics[scale=0.5]{conceptual.png}
\end{center}
\let\thefootnote\relax\footnote{
	\begin{itemize}
		\item Se considera que los asientos se van llenando de forma ordenada, de modo que el número de asiento permite identificar el último asiento posible.
		\item Se relacionó comentarios con usuario, ya que son estos los que comentan las películas, por lo que se deja para la implementación decidir si se hace algo para los empleados que deseen comentar.
	\end{itemize}
}

\newpage
\section{Modelo Lógico}
\vspace{0.6cm}
\begin{center}
	\includegraphics[scale=0.5]{logico.png}
\end{center}
\footnote{
	\begin{itemize}
		\item En actores y directores se tomó el nombre como PK, ya que al momento de que el empleado ingrese uno nuevo, no hay forma de determinar si estaba ya ingresado más que con dicho nombre.
		\item Para facilitar el trabajo se crearon algunas claves primarias como: id$\_$funcion o id$\_$cliente.
	\end{itemize}
}
\end{document}
